% 问题四专用草稿，整合进总论文前请与问题一至三的评价合并。
\section{模型评价}

\subsection{模型的优点}
\begin{itemize}[label=$\checkmark$]
  \item 问题四第二问将波动电价纳入两阶段随机规划，结合期望费用与条件风险价值，权衡平均成本和高费用风险。
  \item 从已结束的历史日联合抽取价格、负荷与光伏残差，保留预测误差的时间相关性及变量间联系。
  \item 问题四第三问利用新发布的光伏预报调整未执行时段，回测表明六点和十二点更新能够降低实际费用。
\end{itemize}

\subsection{模型缺点及改进}
\begin{itemize}[label=\textbullet]
  \item 历史残差难以覆盖未出现的极端波动，可补充有数据依据的压力场景。
  \item 日内价格预测未随已实现价格更新，固定时点重算的收益也不总是明显，可引入因果价格更新和收益触发机制。
  \item 模型允许计划购电量未全部提取，推广前需核实合同规则，并据此调整执行约束。
\end{itemize}
